\section{Related Work}
\subsection{Train dLLMs from scratch}
Auto-regressive language models~\citep{lingv2,moonshot2025kimik2,liu2024deepseekv3,meta2025llama4} are typically trained by maximizing the likelihood of predicting the next token. Under this paradigm, model performance has been shown to scale effectively with increasing model size, dataset volume, and computational resources, following well-established scaling laws. Recently, MDLMs~\citep{song2025seed,Dream7B2025,nie2025llada} have emerged as an alternative generative framework, reformulating text generation as an iterative denoising process. In each forward step, a subset of tokens is randomly masked, and the model is trained to recover the original tokens conditioned on the remaining unmasked context. %

Encouraged by this paradigm shift, several studies have explored training MDLMs from scratch to assess their full potential. For instance, LLaDA~\citep{nie2025llada} demonstrated that a 8B dense MDLM, trained entirely from scratch, achieves performance competitive with similarly sized AR counterparts. Building upon this, LLaDA-MoE~\citep{zhu2025lladamoe} introduced the Mixture-of-Experts (MoE) architecture into the MDLM for the first time, showing that a scratch-trained MoE-based MDLM can surpass dense models in both efficiency and capability, thereby validating the compatibility and scalability of MDLMs with advanced MoE designs. Moreover, due to the fundamentally different training dynamics compared to AR models, established training practices and hyperparameter recipes from the AR domain are often suboptimal for MDLMs. To address this gap, recent efforts such as Quakka~\citep{quakka2025} and OpenMoE2~\citep{openmoe2} have begun investigating the scaling properties and optimal training strategies specifically tailored for MDLMs, laying the groundwork for principled scaling in this emerging paradigm.

However, from-scratch trained MDLMs still lag behind state-of-the-art AR models in overall performance. This gap can be largely attributed to the disparity in training data volume and the maturity of infrastructure support—factors that have been extensively optimized over years of development for AR models. Moreover, due to the high computational cost and long training cycles required for pretraining from scratch, MDLMs mentioned above are typically limited in model scale ($\leq$8B), whereas leading AR models now routinely scale into tens or even hundreds of billions. 

\subsection{Scaling dLLMs with AR initialization}
Given the strong knowledge capacity and performance of AR models, several recent studies have explored initializing dLLMs from pre-trained AR models to reduce training costs and narrow the performance gap between AR models and dLLMs. For instance, DiffusionLLaMA~\citep{diffullama} and Dream-7B~\citep{Dream7B2025} adopt a mask annealing strategy to gradually transition from causal attention to bidirectional attention during training, while employing a CART-based loss reweighting scheme to balance token-level learning dynamics. In contrast, RND1~\citep{RND1} takes a more direct approach by immediately converting the causal attention mechanism of the AR model into a bidirectional one upon initialization. Notably, RND1 observes that when initializing DLM training from an AR model, preserving knowledge-intensive capabilities requires constraining updates to the model’s dense layers to prevent catastrophic forgetting.

Block Diffusion Language Models (BDLMs)~\citep{Blockdiffusion2025} provide a hybrid paradigm that balances efficiency and performance by combining diffusion and AR modeling. Tokens are generated block-wise: within each block, a diffusion process reconstructs masked tokens, while blocks are produced auto-regressively. 
This design enables variable-length generation and supports KV-cache reuse during decoding, enhancing inference efficiency. Consequently, BDLMs can be effectively initialized from AR models, narrowing the performance gap. For example, SDAR~\citep{SDAR2025} leverages the Qwen-3 series~\citep{yang2025qwen3} to train more efficient BDLMs. By exploring various block sizes and optimization strategies, it achieves performance comparable to its AR base model. 

However, one key limitation across all existing methods is their restricted model scale—ranging only from 7B to 30B parameters—leaving the feasibility and scalability of AR-initialized diffusion models largely unexplored at larger scales. 
Besides, the low training efficiency of block diffusion hinders its widely application to large-scale corpus for large-size models. 
Whether such initialization strategies can effectively generalize to models beyond the 30B scale remains an open question.


\subsection{dLLMs post-training}
Beyond pre-training, post-training is crucial for unlocking the full potential of dLLMs by aligning them with specific tasks and human preferences. This process typically involves supervised fine-tuning (SFT) to instill instruction-following capabilities, reinforcement learning (RL) to enhance complex reasoning, and inference optimization to address efficiency bottlenecks.

Recent work has explored SFT to adapt dLLMs for specialized domains. For instance, Dream-Coder~\citep{xie_dream-coder_2025} fine-tunes a 7B dLLM for code generation, demonstrating unique abilities like adaptive "sketch-then-fill" strategies for complex algorithms. Similarly, the general-purpose model Dream-7B~\citep{Dream7B2025} leverages SFT to achieve performance on par with top-tier AR models, while uniquely excelling at tasks requiring complex planning and constraint satisfaction. Other studies have investigated specialized fine-tuning strategies to balance quality and efficiency. Seed-Diffusion~\citep{song2025seed}, for example, employs a two-stage curriculum learning strategy to train a high-speed code generation model, while LiDAR~\citep{liu_tidar_2025} introduces a hybrid "think in diffusion, generate in AR" architecture through fine-tuning, significantly boosting inference throughput while maintaining quality.

To further enhance dLLMs' reasoning abilities, researchers have begun adapting reinforcement learning techniques. However, applying standard policy gradient methods is challenging due to the intractable log-likelihood of dLLMs. To address this, SPG~\citep{wang2025spg} proposes a novel Sandwich Policy Gradient algorithm that obtains a more robust and less biased gradient by maximizing an evidence lower bound for high-reward samples and minimizing an evidence upper bound for low-reward ones. Another line of work, TraceRL~\citep{wang2025revolutionizing}, focuses on aligning the training objective with the model's multi-step generation trajectory. This framework led to the TraDo series of models, which have not only surpassed strong AR models on reasoning benchmarks but also produced the first dLLM capable of long-chain-of-thought reasoning.

A significant challenge for dLLMs is their slow inference speed, stemming from the iterative nature of the denoising process. To mitigate this, several acceleration methods have been proposed. DPad~\citep{chen_dpad_2025} offers a training-free solution by treating future tokens as a dynamic "scratchpad" and using a sliding window and distance-based pruning to reduce redundant computations, achieving a dramatic speedup, especially for long sequence generation. In contrast, D2F~\citep{wang_diffusion_2025} introduces a hybrid autoregressive-diffusion paradigm that enables parallel denoising of future text blocks even before preceding ones are fully generated. This approach allows dLLMs to leverage KV-caching and, for the first time, surpass the inference speed of equivalently sized AR models.

Despite these advances, the field of dLLM post-training is still nascent. Systematic exploration of how these techniques—SFT, RL, and acceleration—interact with one another, and how they scale to models with hundreds of billions of parameters, remains an open and critical area for future research.
